\documentclass{beamer}
\mode<presentation> {
%\usetheme{Dresden}
\usetheme{CambridgeUS}
}
\newcommand{\subitem}{\par\qquad}
\usepackage{booktabs}

\usepackage{tikz}
\usepackage{subfigure}	
\usepackage[round]{natbib}   % omit 'round' option if you prefer square brackets
\DeclareMathOperator*{\argmin}{argmin}
\DeclareMathOperator*{\argmax}{argmax}
\newtheorem{lem}{Lemma}
\newtheorem{thm}{Theorem}
\newtheorem{nots}{Note}
\newtheorem{defs}{Definition}
\usetikzlibrary{arrows,calc,positioning,decorations.pathreplacing} 
\tikzset{
	>=stealth',
	punkt/.style={
		circle,
		draw=black, thick,
		minimum height=2em,
		inner sep=0pt,
		text centered},
	pil/.style={
		->,
		thick,
		shorten <=2pt,
		shorten >=2pt,},
	cir/.style={
		rectangle,
		rounded corners,
		draw=black, thick,
		minimum width=1.5em,
		text centered},
	arr/.style={
		->,
		thick,
		shorten <=2pt,
		shorten >=2pt,},
	layer/.style={
		draw=black,
		fill=white,
		thick,
		rectangle,
		rounded corners,
		minimum width=1.5em,
		minimum height=12em
	},
	brace/.style={
		thick,
		decoration={
			brace,
			mirror,
			raise=2.5cm
		},
		decorate
	}
}

\newcommand{\drawlayer}[4]{
	\node[layer, draw=#4] (#1) at (#2,#3) {};
	\foreach \x in {-1.75, -1.25, -0.75, -0.25, 0.25, 0.75, 1.25, 1.75} \draw[thick, #4] (#2,\x+#3) circle (0.15);
}
\newcommand\irregularcircle[2]{
	\pgfextra {\pgfmathsetmacro\len{(#1)+rand*(#2)}}
	+(0:\len pt)
	\foreach \a in {10,20,...,350}{
		\pgfextra {\pgfmathsetmacro\len{(#1)+rand*(#2)}}
		-- +(\a:\len pt)
	} -- cycle
}

\newcommand{\independent}{\mbox{${}\perp\mkern-11mu\perp{}$}}
\newcommand{\notindependent}{\mbox{${}\not\!\perp\mkern-11mu\perp{}$}}
\newcommand{\iid}{\overset{\text{iid}}{\sim}}

%%%%%%%%%%%%%%%%%%%%%%%%%%%%%%%%%%%%%%%%%%%%%%%%%%%%%%%%%%%%%%%%%%%%%%%%%%
% PARENTHESES 
%%%%%%%%%%%%%%%%%%%%%%%%%%%%%%%%%%%%%%%%%%%%%%%%%%%%%%%%%%%%%%%%%%%%%%%%%%

\renewcommand{\dot}[2]{\langle {#1},\, {#2} \rangle}
\newcommand{\pa}[1]{\left( #1 \right)}
\newcommand{\pb}[1]{\left\{ #1 \right\}}
\newcommand{\pc}[1]{\left[ #1 \right]}
\newcommand{\pn}[1]{\left\| #1 \right\|}
\newcommand{\ps}[1]{\left| #1 \right|}
\newcommand{\ot}{\leftarrow}
\newcommand{\given}{\, | \,}

%%%%%%%%%%%%%%%%%%%%%%%%%%%%%%%%%%%%%%%%%%%%%%%%%%%%%%%%%%%%%%%%%%%%%%%%%%
% MATHCAL
%%%%%%%%%%%%%%%%%%%%%%%%%%%%%%%%%%%%%%%%%%%%%%%%%%%%%%%%%%%%%%%%%%%%%%%%%%

\newcommand{\D}{\mathcal{D}}
\renewcommand{\P}{\mathcal{P}}
\renewcommand{\Pr}{\mathbb{P}}
\renewcommand{\H}{\mathcal{H}}
\renewcommand{\L}{\mathcal{L}}
\newcommand{\Hk}{{\mathcal{H}_k}}
\newcommand{\F}{\mathcal{F}}
\newcommand{\X}{\mathcal{X}}
\newcommand{\Y}{\mathcal{Y}}
\newcommand{\Z}{\mathcal{Z}}
\newcommand{\N}{\mathcal{N}}
\newcommand{\Ex}{\mathcal{E}}
\newcommand{\Hyp}{\mathcal{H}}
\newcommand{\B}{\mathcal{B}}


%%%%%%%%%%%%%%%%%%%%%%%%%%%%%%%%%%%%%%%%%%%%%%%%%%%%%%%%%%%%%%%%%%%%%%%%%%
% MATHBB 
%%%%%%%%%%%%%%%%%%%%%%%%%%%%%%%%%%%%%%%%%%%%%%%%%%%%%%%%%%%%%%%%%%%%%%%%%%

\newcommand{\Pio}{\mathbb{P}}
\newcommand{\R}{\mathbb{R}}
\newcommand{\Rd}{\mathbb{R}^d}
\newcommand{\Rq}{\mathbb{R}^q}
\newcommand{\Rn}{\mathbb{R}^n}
\newcommand{\Rm}{\mathbb{R}^m}
\newcommand{\Rnd}{\mathbb{R}^{n \times d}}
\newcommand{\Rnm}{\mathbb{R}^{n \times m}}
\newcommand{\I}{\mathbb{I}}
\newcommand{\Pa}{\mathrm{Pa}}
\renewcommand{\d}{\mathrm{d}}
\newcommand{\sig}{\mathrm{sign}\circ\!}
\newcommand{\Rp}{R_{\varphi}}
\newcommand{\Rpn}{\hat{R}_{\varphi}}
\newcommand{\Rpnt}{\tilde{R}_{\varphi}}
\newcommand{\f}{f^*}
\newcommand{\fn}{\hat{f}_n}
\newcommand{\fnt}{\tilde{f}_n}
\newcommand{\fntm}{\tilde{f}^m_n}
\newcommand{\Mk}{\M_k}
\newcommand{\muP}{\mu_k(\P)}
\renewcommand{\d}{{\mathrm{d}}}
\newcommand{\Rad}{\mathrm{Rad}}
\newcommand{\M}{\mathscr{M}}
\newcommand{\head}[2]{\multicolumn{1}{>{\centering\arraybackslash}p{#1}}{\textbf{#2}}}
\renewcommand{\do}{do}
\title[de Finetti Theorem]{Subjective Theory of Probability} 
\subtitle{Dutch Book (de Finetti's) Theorem}

\author[Behrad Moniri]{Behrad Moniri}

\institute[]
{

La soci\'{e}t\'{e} secr\'{e}te des statisticiens\\ et des ing\'{e}nieurs \'{e}lectriciens au ch\^{o}mage\\\vspace{0.5cm}
\texttt{bemoniri@live.com}
}
\date{}

%	\includegraphics[width=2cm]{digilogo.png}


\newcommand{\bigCI}{\mathrel{\text{\scalebox{1.07}{$\perp\mkern-10mu\perp$}}}}

\begin{document}
\maketitle

\begin{frame}
\frametitle{Introduction}
\begin{itemize}

\item The most widely accepted view interprets probabilities:
long run averages. This is based on the fact that averages should settle down to expectations over a long sequence of \textit{independent} trials.


\pause
\item de Finetti theorem provides an alternative view that does not depend on a preliminary concept of independence, and which concentrates attention on the \textit{linearity properties} of expectations.

\end{itemize}
\end{frame}

\begin{frame}
\begin{itemize}
\item
Let $\Omega$ be your sample space and $X: \Omega \to \mathbb{R}$ be a bounded function (random variable). 
\pause
\item
Forget about probability measures on $\Omega$. Suppose you consider $p(X)$ to be the \textit{fair price}  to pay now in order to receive $X$ at some later time.
\pause
\item By fair I mean that you should be prepared to accept a payment $p(X)$ from me now in return for giving me an amount $X$ later.
\pause
\item Your return: $X'(\omega) = X(\omega) - p(X)$. We call this \textit{fair return}.
\end{itemize}
\end{frame}


\begin{frame}
\frametitle{Properties of fair bets!}

Unless you start worrying about utilities, you should find the following properties reasonable:
\pause
\begin{itemize}
	\item
	\textbf{fair + fair = fair}. That is, if you consider $p(X)$ fair for $X$ and $p(Y)$ fair for $Y$, then you should be prepared to make both bets, paying $p(X) + p(Y)$ to receive $X + Y$.
	\pause
	\item
	\textbf{constant $\times$ fair = fair}. You shouldn't object if I suggest you pay $cp(X)$ to receive $cX$ for constant $c$.

\end{itemize}

\pause

These two conditions imply that imply that the collection of all fair returns is a vector space over field $\mathbb{R}$.
\end{frame}

\begin{frame}
\frametitle{Properties of fair bets!}
There is a third reasonable property that goes by several names: \textit{coherency} or
\textit{nonexistence of a Dutch book}, the \textit{no-arbitrage requirement}, or \textit{the no-free-lunch principle}:
\pause
\begin{itemize}
	\item
	There is no fair return $X'$ for which $X'(\omega) < 0$ for all $\omega \in \Omega$, with strict inequality for at least one $\omega$.
\end{itemize}
\end{frame}


\begin{frame}
\frametitle{Properties of fair bets!}
\begin{lem}
The previous properties imply that $p(.)$ is a linear, increasing functional on random variables.
\end{lem}

\begin{proof}
For constants $\alpha$ and $\beta$ and random variables $X$ and $Y$ with fair prices $p(X)$ and $p(Y)$, consider the combined effect of the following fair bets:
\begin{itemize}
\item
You pay me $\alpha p(X)$ to receive $\alpha X$.
\item
You pay me $\beta p(Y)$ to receive $\beta Y$.
\item
I pay you $p(\alpha X + \beta Y)$ to receive $\alpha X + \beta Y$.
\end{itemize}

Your net return is $c = p(\alpha X + \beta Y) - \alpha p(X) - \beta p(Y)$. 

If $c>0$, (iii) is violated. If $c<0$, consider the other side bet to violate (iii). This proves linearity.

\end{proof}
\end{frame}

\begin{frame}
\begin{proof}

To prove that $p(.)$ is increasing, suppose $\forall \omega \in \Omega:\;  X(\omega) \geq Y(\omega)$.

 If you claim
that $p(X) < p(Y)$ then I would be happy for you to accept the bet that delivers $(Y - p(Y)) - (X - p(X)) =
-(X - Y) - (p(Y) - p(X))$,
which is always $< 0$.
\end{proof}

\begin{nots}
If both $X$ and $X-p(X)$ are fair, so is $X - \big(X-p(X)\big)$ with constant return. This imples that $p(X) = 0$.
\end{nots}

\end{frame}

\begin{frame}
\frametitle{de Finettin Theorem}
\begin{thm}
 $	 p(F_X \cup F_Y) = p(F_X) + p(F_Y)$ for disjoint $F_X, F_Y \subseteq \Omega$. Here we have used the de Finetti notation $p(A) = p(\mathbf{1}_A)$ for $A \subseteq \Omega$.
\end{thm}
\pause
\begin{proof}
	As a special case, consider the bet that returns $1$ if an event $F$ occurs, and $0$	otherwise. The previous theorem  implies   $$	 p(F_X \cup F_Y) =	p(F_X) + p(F_Y)$$ for disjoint $F_X$ and $F_Y$.
\end{proof}
\pause
We can similiary show that $p(\Omega) = 1$ and $p(\emptyset) = 0$.

\end{frame}

\begin{frame}
\frametitle{Conditional Probability and Contingent Bets}
Things become much more interesting if you are prepared to make a bet to receive an amount $X$ but only when some event $F$ occurs. 
\pause

\begin{itemize}
\item Typically, knowledge of the occurrence of $F$ should change the fair price, which we could denote by $p(X | F)$. 
\pause
\item The bet that returns $\big(X - p(X | F)\big) F$ is fair. 
\pause
\item The indicator function $F$ ensures that money changes hands only when $F$ occurs.
\end{itemize}
\end{frame}

\begin{frame}
\begin{thm}
If $\Omega$ is partitioned into disjoint events $F_1, \dots, F_k$. and $X$ is a random variable, then $p(X) = \sum_{i = 1}^{k} p(Fi)\, p(X | F_i)$. 
\end{thm}
\begin{proof}
For a single $F_i$, argue by linearity that $$0 = p (X F_i - p(X | F_i)F_i) = p(XF_i) - p(X | F_i)\,p(Fi).$$
Sum over $i$, using linearity again, together with the fact that $X = \sum_i X Fi$, to deduce
that $p(X) =\sum_i p(XF_i) = \sum_i p(F_i)p(X | F_i)$, as asserted.
\end{proof}
\end{frame}

\begin{frame}

\begin{itemize}
	\item Why should we restrict the Lemma to finite partitions?
	\pause
	\item  If we allowed countable partitions we would get the countable additivity property-the key requirement in the theory of measures.
		\pause
	\item  If we accept that assumption, then why not accept that arbitrary combinations of fair events are fair?
		\pause
	\item
	 For uncountably infinite collections we would run into awkward contradictions.
	 	\pause	 
	 \item  For example, suppose $\omega$ is generated from a uniform distribution on $[0, 1)$. Let $X_t = \mathbf{1}_{\omega = t}$.
	 	\pause
	  By symmetry one might expect $p(X_t) = c$ for
	 some constant c that doesn't depend on t.
	 \pause However
	 
	 $$1 = p(1) = p(\sum_{0\leq t \leq 1} X_t) \stackrel{?}{=} \sum_{0\leq t \leq 1} p(X_t) = \begin{cases}
	 0 & c = 0\\
	 \infty & \text{else}
	 \end{cases}$$
	 
\end{itemize}




\end{frame}

\end{document}

